\section{Propiedades de permanencia}
\label{II.3}

\begin{proposition}
\label{II.3.1}
Sea $f\colon X\to Y$ un morfismo, sea $x\in X$ y sea $y=f(x)$. Supóngase que $f$ es liso en~$x$. Para que $\cal{O}_x$ sea reducido (\resp regular, \resp normal), es necesario y suficiente que $\cal{O}_y$ lo sea.
\end{proposition}

En efecto, este enunciado es conocido cuando $X$ es de la forma $Y[t_1,\dots,t_n]$, y se ha demostrado en~(I, \no \Ref{I.9}) para un morfismo étale; el caso general se deduce inmediatamente gracias a la definición~\Ref{II.1.1}.

No detallaremos aquí las demás propiedades de permanencia, que resultan ya de la sola planitud, o del hecho de que $X$ sea localmente cuasifinito y plano sobre un $Y$-preesquema de la forma $Y[t_1,\dots,t_n]$ (o, como
\marginpar{35}
diremos, de que $X$ sea
\ifthenelse{\boolean{orig}}{Cohen--Macauley}{Cohen--Macaulay}\ % corrección ortográfica ya incorporada por la autoridad.
\index{Cohen--Macaulay (esquema de)|hyperpage}%
sobre~$Y$). Señalemos solamente que de este último hecho se sigue que
\begin{equation}
\label{eq:II.3.1}
{\dim\cal{O}_x=\dim\cal{O}_y+n-d,\quad
\prof\cal{O}_x=\prof\cal{O}_y+n-d},
\end{equation}
donde $n$ es la dimensión de la fibra de~$f$ en~$x$, y $d$ el grado de trascendencia de~$\kres(x)$ sobre~$\kres(y)$; de donde (poniendo $\coprof=\dim-\prof$)\footnote{Para estas fórmulas, \cf EGA~IV~6.1 y~6.3.}
\begin{equation}
\label{eq:II.3.2}
{\coprof\cal{O}_x=\coprof\cal{O}_y}.
\end{equation}
De ello se sigue, por ejemplo, que $\cal{O}_x$ es Cohen--Macaulay (\resp no tiene componentes embebidas) si y solo si lo mismo ocurre con~$\cal{O}_y$.
